
\section{\trExponentialDecay}
\label{sec:exponential-decay}

在本节中，我们来证明那个我们已经使用过多次的衰减性质，也就是下面的结果。

\begin{theorem}\label{thm:exp-decay-dichotomy}
如果 $Y$ 是沿一个有限\trEnergy{}的解\trLinearization{}后的\trFloerEquation{}的 $\cC^{2}$ 类解，那么：
\begin{itemize}
\item 或者，当 $s$ 趋于 $\pm\infty$ 时，$\int\|Y(s,t)\|^{2}\,dt$ 趋于无穷；
\item 或者，$Y$ 满足
\[
\|Y(s,t)\|\le Ce^{-\delta|s|}
\]
其中 $\delta$ 与 $C$ 是某些常数（并对每个 $t$ 都成立）。
\end{itemize}
\end{theorem}

\begin{remark}\label{thm:du-ds-decay-remark}
如果 $u$ 是我们沿之做\trLinearization{}的那条\trFloerEquation{}的解，那么解 $Y=\partial u/\partial s$ 便具有\trExponentialDecay{}性质，特别地，对每个 $p>1$ 它都属于空间 $L^{p}(\RR\times S^{1};\RR^{2n})$。的确，由于 $u$ 的\trEnergy{}是有限的，
\[
\int\Bigl\|\frac{\partial u}{\partial s}\Bigr\|^{2}\,dt\quad\text{在 }s\longrightarrow\pm\infty\text{ 时不趋于 }+\infty\text{，}
\]
这就为这一类型的解排除了定理陈述中的第一种情形。还应注意，这显然蕴含 $\scr{M}(x,y)\subset\cC^{\infty}_{\searrow}(x,y)$，正如命题~8.2.3~所断言的那样。
\end{remark}

我们将把这条定理证明成下面这条命题的一个推论，该命题断言了一个“\trExponentialDecay”性质。

\begin{proposition}\label{thm:l2-exp-decay}
设 $u$ 是\trFloerEquation{}的一个有限\trEnergy{}的解。如果 $Y$ 是沿 $u$ \trLinearization{}后的\trFloerEquation{}的一个解，使得 $\int\|Y(s,t)\|^{2}\,dt$ 不趋于无穷，那么就存在常数 $\delta>0$（它不依赖于 $u$）与 $C>0$，使得
\[
\int_{0}^{1}\|Y(s,t)\|^{2}\,dt\le Ce^{-\delta|s|}\text{。}
\]
\end{proposition}

\begin{proof}
公式
\[
f(s)=\frac{1}{2}\int_{0}^{1}\|Y(s,t)\|^{2}\,dt=\frac{1}{2}\,\|Y\|^{2}_{L^{2}(S^{1})}
\]
定义了一个 $\cC^{2}$ 类的函数 $f\colon\RR\to\RR$。我们将要证明这个函数满足一条\trDifferential{}方程
\[
f''\ge\delta^{2}f
\]
其中 $\delta$ 是某个常数。于是让我们来计算它的各阶导数：
\[
\begin{gathered}
f'(s)=\int_{0}^{1}\Bigl\langle Y(s,t),\frac{\partial Y}{\partial s}(s,t)\Bigr\rangle\,dt\\
f''(s)=\int_{0}^{1}\Bigl\|\frac{\partial Y}{\partial s}(s,t)\Bigr\|^{2}\,dt+\int_{0}^{1}\Bigl\langle Y(s,t),\frac{\partial^{2}Y}{\partial s^{2}}(s,t)\Bigr\rangle\,dt\text{。}
\end{gathered}
\]
我们利用 $Y$ 是\trLinearization{}后的方程的一个解——也就是说它满足
\[
\frac{\partial Y}{\partial s}=-J_{0}\frac{\partial Y}{\partial t}-SY
\]
——来确定第二项。于是我们有
\[
\frac{\partial^{2}Y}{\partial s^{2}}=-J_{0}\frac{\partial^{2}Y}{\partial s\,\partial t}-\frac{\partial S}{\partial s}Y-S\frac{\partial Y}{\partial s}\text{。}
\]
因此
\[
\int_{0}^{1}\Bigl\langle Y,\frac{\partial^{2}Y}{\partial s^{2}}\Bigr\rangle\,dt=-\int_{0}^{1}\Bigl\langle Y,J_{0}\frac{\partial^{2}Y}{\partial s\,\partial t}\Bigr\rangle\,dt-\int_{0}^{1}\Bigl\langle Y,\frac{\partial S}{\partial s}Y\Bigr\rangle\,dt-\int_{0}^{1}\Bigl\langle Y,S\frac{\partial Y}{\partial s}\Bigr\rangle\,dt\text{。}
\]
这些积分中的第一个可以用分部积分来计算：
\[
\begin{aligned}
\int_{0}^{1}\Bigl\langle Y,J_{0}\frac{\partial^{2}Y}{\partial s\,\partial t}\Bigr\rangle\,dt
&=-\int_{0}^{1}\Bigl\langle\frac{\partial Y}{\partial t},J_{0}\frac{\partial Y}{\partial s}\Bigr\rangle\,dt\\
&=-\int_{0}^{1}\Bigl\langle J_{0}\frac{\partial Y}{\partial s},J_{0}\frac{\partial Y}{\partial s}\Bigr\rangle\,dt-\int_{0}^{1}\Bigl\langle J_{0}SY,J_{0}\frac{\partial Y}{\partial s}\Bigr\rangle\,dt\\
&=-\int_{0}^{1}\Bigl\|\frac{\partial Y}{\partial s}\Bigr\|^{2}\,dt-\int_{0}^{1}\Bigl\langle SY,\frac{\partial Y}{\partial s}\Bigr\rangle\,dt\text{。}
\end{aligned}
\]
于是我们得到
\[
\begin{aligned}
\int_{0}^{1}\Bigl\langle Y,\frac{\partial^{2}Y}{\partial s^{2}}\Bigr\rangle\,dt={}&\int_{0}^{1}\Bigl\|\frac{\partial Y}{\partial s}\Bigr\|^{2}\,dt+\int_{0}^{1}\Bigl\langle SY,\frac{\partial Y}{\partial s}\Bigr\rangle\,dt\\
&-\int_{0}^{1}\Bigl\langle Y,\frac{\partial S}{\partial s}Y\Bigr\rangle\,dt-\int_{0}^{1}\Bigl\langle Y,S\frac{\partial Y}{\partial s}\Bigr\rangle\,dt\text{。}
\end{aligned}
\]
函数 $f$ 的二阶导数因而为\footnote{译注：原书此式第二项印作 $\int_{0}^{1}\bigl\langle Y,\frac{\partial S}{\partial s}\bigr\rangle\,dt$；按本页上方诸式的直接计算以及下文（p.298）估计中所用的 $\int_{0}^{1}\bigl\langle Y,\frac{\partial S}{\partial s}Y\bigr\rangle\,dt$，疑为漏印因子 $Y$，此处径改。又，首项按直算应为 $2\int_{0}^{1}\bigl\|\frac{\partial Y}{\partial s}\bigr\|^{2}\,dt$（这一出入不影响下文的不等式链与其结论），此处照原书排印，存疑备案。——译者注}
\[
f''(s)=\int_{0}^{1}\Bigl\|\frac{\partial Y}{\partial s}\Bigr\|^{2}\,dt-\int_{0}^{1}\Bigl\langle Y,\frac{\partial S}{\partial s}Y\Bigr\rangle\,dt+\int_{0}^{1}\Bigl\langle Y,({}^{t}S-S)\frac{\partial Y}{\partial s}\Bigr\rangle\,dt\text{。}
\]
利用这一事实
\[
\lim_{s\to\pm\infty}\ \sup_{t}\ \frac{\partial S}{\partial s}(s,t)=0
\]
（它已在\cref{thm:df-dbar-plus-s}中获得证明）以及\trCauchy--\trSchwarz{}不等式，我们便得到
\[
\begin{aligned}
\int_{0}^{1}\Bigl\langle Y,\frac{\partial S}{\partial s}Y\Bigr\rangle\,dt
&\le\Bigl(\int_{0}^{1}\|Y\|^{2}\,dt\Bigr)^{1/2}\Bigl(\int_{0}^{1}\Bigl\|\frac{\partial S}{\partial s}Y\Bigr\|^{2}\,dt\Bigr)^{1/2}\\
&\le\Bigl(\int_{0}^{1}\|Y\|^{2}\,dt\Bigr)^{1/2}\varepsilon\Bigl(\int_{0}^{1}\|Y\|^{2}\,dt\Bigr)^{1/2}\\
&\le\varepsilon\int_{0}^{1}\|Y\|^{2}\,dt\qquad\text{对充分大的 }|s|\text{。}
\end{aligned}
\]
由于 ${}^{t}S-S$ 趋于 $0$，我们有
\[
\Bigl|\int_{0}^{1}\Bigl\langle Y,({}^{t}S-S)\frac{\partial Y}{\partial s}\Bigr\rangle\,dt\Bigr|\le\eta\Bigl(\int_{0}^{1}\|Y\|^{2}\,dt\Bigr)^{1/2}\Bigl(\int_{0}^{1}\Bigl\|\frac{\partial Y}{\partial s}\Bigr\|^{2}\,dt\Bigr)^{1/2}
\]
对充分大的 $|s|$ 成立。于是我们便得到不等式
\[
f''(s)\ge\int_{0}^{1}\Bigl\|\frac{\partial Y}{\partial s}\Bigr\|^{2}\,dt-\varepsilon\int_{0}^{1}\|Y\|^{2}\,dt-\eta\Bigl(\int_{0}^{1}\|Y\|^{2}\,dt\Bigr)^{1/2}\Bigl(\int_{0}^{1}\Bigl\|\frac{\partial Y}{\partial s}\Bigr\|^{2}\,dt\Bigr)^{1/2}\text{。}
\]
现在让我们来利用这一事实\footnote{译注：原书此处末一等号印作“$=AY$”；按 $AY=A_{0}Y+(S-S^{\pm})Y$（即 $A=J_{0}\frac{\partial}{\partial t}+S$）实应有 $AY=-\frac{\partial Y}{\partial s}$，疑为漏印负号，此处径改（下文仅用到 $\|AY\|$，不受影响）。又，其下 $f''(s)$ 等式链首行原书亦印作 $=$，按上文 $f''$ 的估计应作 $\ge$，兹照原书排印，存疑备案。——译者注}
\[
\frac{\partial Y}{\partial s}=-J_{0}\frac{\partial Y}{\partial t}-SY=-AY\text{，}\qquad\text{其中 }AY=A_{0}Y+(S-S^{\pm})Y\text{，}
\]
其中 $A_{0}$ 是 $L^{2}$ 上的一个算子，我们已在\cref{thm:a-invertible-p2}中证明它是\trInvertible{}的，而 $S-S^{\pm}$ 在 $s\to\pm\infty$ 时趋于 $0$。因此，对充分大的 $|s|$，$A$ 是\trInvertible{}的（视为从 $W^{1,2}(S^{1})$ 到 $L^{2}(S^{1})$ 的算子），并且我们有
\[
\|AY\|_{L^{2}(S^{1})}\ge D\,\|Y\|_{W^{1,2}(S^{1})}\ge D\,\|Y\|_{L^{2}(S^{1})}
\]
对某个正常数 $D$。于是我们便有
\[
\begin{aligned}
f''(s)&=\|AY\|^{2}_{L^{2}}-\varepsilon\|Y\|^{2}_{L^{2}}-\eta\,\|Y\|_{L^{2}}\|AY\|_{L^{2}}\\
&=\|AY\|_{L^{2}}\left(\|AY\|_{L^{2}}-\eta\,\|Y\|_{L^{2}}\right)-\varepsilon\|Y\|^{2}_{L^{2}}\\
&\ge\left(D(D-\eta)-\varepsilon\right)\|Y\|^{2}_{L^{2}}\\
&\ge\delta^{2}f(s)
\end{aligned}
\]
对充分大的 $|s|$ 与某个常数 $\delta>0$ 成立。函数 $f$ 因而满足一条所论类型的\trDifferential{}方程。该命题证明的收尾来自下面这条引理（我们在 $+\infty$ 处陈述它，但它在 $-\infty$ 处的类似结论显然也成立）。

\begin{lemma}\label{thm:f-growth-dichotomy}
如果 $f''(s)\ge\delta^{2}f(s)$ 对 $s\ge s_{0}$ 成立，那么：
\begin{itemize}
\item 或者 $\lim_{s\to+\infty}f(s)=+\infty$；
\item 或者 $f(s)\le f(s_{0})e^{-\delta(s-s_{0})}$ 对 $s\ge s_{0}$ 成立。
\end{itemize}
\end{lemma}

\begin{proof}
令
\[
g(s)=e^{-\delta s}\left(f'(s)+\delta f(s)\right)\text{，}
\]
于是
\[
\begin{aligned}
g'(s)&=e^{-\delta s}\left(-\delta(f'(s)+\delta f(s))+f''(s)+\delta f'(s)\right)\\
&=e^{-\delta s}\left(f''(s)-\delta^{2}f(s)\right)\\
&\ge 0\text{，}\qquad\text{由对 }f\text{ 的假设。}
\end{aligned}
\]
特别地，函数 $g$ 在 $[s_{0},+\infty[$ 上是递增的。

让我们首先假设存在一个 $s_{1}\ge s_{0}$ 使得 $g(s_{1})>0$。那么对 $s\ge s_{1}$ 就有 $g(s)\ge g(s_{1})$，从而
\[
f'(s)+\delta f(s)\ge e^{\delta s}g(s_{1})
\]
对 $s\ge s_{1}$ 成立。令 $h(s)=e^{\delta s}f(s)$，那么
\[
h'(s)=e^{\delta s}(\delta f+f')\ge e^{2\delta s}g(s_{1})\text{，}
\]
于是 $h(s)\ge Ae^{2\delta s}$ 对 $s\ge s_{1}$ 成立，从而 $f(s)\ge Ae^{\delta s}$ 趋于 $+\infty$。

如果这样的 $s_{1}$ 不存在，那么对 $s\ge s_{0}$ 我们有 $g(s)\le 0$，从而也有 $f'(s)+\delta f(s)\le 0$，于是 $h(s)=e^{\delta s}f(s)$ 在 $[s_{0},+\infty[$ 上是递减的。此外，
\[
e^{\delta s}f(s)=h(s)\le h(s_{0})=e^{\delta s_{0}}f(s_{0})\text{，}
\]
也就是说，
\[
f(s)\le e^{\delta s_{0}}f(s_{0})e^{-\delta s}\text{，}
\]
这就结束了引理的证明。
\end{proof}

而这也就结束了\cref{thm:l2-exp-decay}的证明。
\end{proof}

我们现在已经证明了
\[
\forall s\in\RR\text{，}\quad\ \|Y\|^{2}_{L^{2}(S^{1})}\le e^{-\delta|s|}\text{，}
\]
但这还不完全是我们所致力的那条定理的陈述。我们还需要验证（逐点的）
\[
\forall(s,t)\in\RR\times S^{1}\text{，}\quad|Y(s,t)|\le Ce^{-\delta|s|}\text{。}
\]
我们将需要下面几个结果。

\begin{lemma}\label{thm:laplacian-lower-bound}
设 $Y$ 是\trLinearization{}后的\trFloerEquation{}的一个 $\cC^{2}$ 类解。那么存在一个常数 $a>1$，使得
\[
\Delta\|Y\|^{2}\ge-a\|Y\|^{2}\text{。}
\]
\end{lemma}

\begin{proposition}\label{thm:mean-value-inequality}
设 $w\colon\RR^{2}\to\RR$ 是一个 $\cC^{2}$ 类的正函数，满足 $\Delta w\ge-aw$，其中 $a>1$ 是一个常数。那么我们有
\[
\forall(s_{0},t_{0})\in\RR^{2}\text{，}\quad w(s_{0},t_{0})\le\frac{8a}{\pi}\int_{B_{1}(s_{0},t_{0})}w\,ds\,dt\text{。}
\]
\end{proposition}

现在，我们只需把这个命题应用于 $w(s,t)=\|Y(s,t)\|^{2}$。

\begin{proof}[\cref{thm:laplacian-lower-bound}的证明]
这是一个简短的计算：
\[
\Delta\langle Y,Y\rangle=\frac{\partial^{2}}{\partial t^{2}}\langle Y,Y\rangle+\frac{\partial^{2}}{\partial s^{2}}\langle Y,Y\rangle=2\Bigl[\Bigl\|\frac{\partial Y}{\partial s}\Bigr\|^{2}+\Bigl\|\frac{\partial Y}{\partial t}\Bigr\|^{2}+\langle\Delta Y,Y\rangle\Bigr]
\]
以及
\[
\Delta=\Bigl(\frac{\partial}{\partial s}-J_{0}\frac{\partial}{\partial t}\Bigr)\Bigl(\frac{\partial}{\partial s}+J_{0}\frac{\partial}{\partial t}\Bigr)\text{，}
\]
于是，若 $Y$ 是\trLinearization{}后的\trFloerEquation{}的一个解，则
\[
\Delta Y=\Bigl(\frac{\partial}{\partial s}-J_{0}\frac{\partial}{\partial t}\Bigr)\Bigl(\frac{\partial}{\partial s}+J_{0}\frac{\partial}{\partial t}\Bigr)Y=\Bigl(\frac{\partial}{\partial s}-J_{0}\frac{\partial}{\partial t}\Bigr)(-SY)\text{。}
\]
因此
\[
\langle\Delta Y,Y\rangle=\Bigl\langle-\frac{\partial S}{\partial s}Y,Y\Bigr\rangle-\Bigl\langle S\frac{\partial Y}{\partial s},Y\Bigr\rangle+\Bigl\langle J_{0}\frac{\partial S}{\partial t}Y,Y\Bigr\rangle+\Bigl\langle J_{0}S\frac{\partial Y}{\partial t},Y\Bigr\rangle
\]
以及
\[
\begin{aligned}
\Delta\langle Y,Y\rangle={}&2\Bigl(\Bigl\|\frac{\partial Y}{\partial s}\Bigr\|^{2}-\Bigl\langle S\frac{\partial Y}{\partial s},Y\Bigr\rangle\Bigr)+2\Bigl(\Bigl\|\frac{\partial Y}{\partial t}\Bigr\|^{2}+\Bigl\langle J_{0}S\frac{\partial Y}{\partial t},Y\Bigr\rangle\Bigr)\\
&-2\Bigl\langle\frac{\partial S}{\partial s}Y,Y\Bigr\rangle+2\Bigl\langle J_{0}\frac{\partial S}{\partial t}Y,Y\Bigr\rangle\text{。}
\end{aligned}
\]
右边的每一项都有一个形如 $-c\|Y\|^{2}$ 的下界。对最后两项而言，这是显然的（由\trCauchy--\trSchwarz{}不等式）。我们于是来研究前两项：\footnote{译注：原书此处首行第二项印作 $\bigl\langle Y\frac{\partial S}{\partial s},Y\bigr\rangle$；按上行 $\Delta\langle Y,Y\rangle$ 表达式首个括号中的 $\bigl\langle S\frac{\partial Y}{\partial s},Y\bigr\rangle$（它与次行的 $\bigl\langle\frac{\partial Y}{\partial s},{}^{t}SY\bigr\rangle$ 配平），疑为 $Y$ 与 $\frac{\partial S}{\partial s}$ 误倒，此处径改。——译者注}
\[
\begin{aligned}
\Bigl\|\frac{\partial Y}{\partial s}\Bigr\|^{2}-\Bigl\langle S\frac{\partial Y}{\partial s},Y\Bigr\rangle
&=\Bigl\|\frac{\partial Y}{\partial s}\Bigr\|^{2}-\Bigl\langle\frac{\partial Y}{\partial s},{}^{t}SY\Bigr\rangle\\
&=\Bigl\|\frac{\partial Y}{\partial s}-\frac{1}{2}{}^{t}SY\Bigr\|^{2}-\frac{1}{4}\|{}^{t}SY\|^{2}\\
&\ge-\frac{1}{4}\|{}^{t}SY\|^{2}\\
&\ge-c\|Y\|^{2}
\end{aligned}
\]
对第二项也有同样的结论，这就给出了所期待的下界。
\end{proof}

\begin{proof}[\cref{thm:mean-value-inequality}的证明]
它从一个引理开始。

\begin{lemma}\label{thm:mean-inequality-ball}
设 $w\colon B_{r}(s_{0},t_{0})\to\RR$ 是一个 $\cC^{2}$ 类的正函数，满足 $\Delta w\ge-b$，其中 $b>0$ 是一个常数。那么我们有\footnote{译注：原书此处第二项印作 $\frac{1}{r^{2}}\int_{B_{r}(s_{0},t_{0})}w$；按本引理证明所得以及下文（p.302）的应用均带 $\pi$，疑为漏印 $\pi$，此处径改。——译者注}
\[
w(s_{0},t_{0})\le\frac{br^{2}}{8}+\frac{1}{\pi r^{2}}\int_{B_{r}(s_{0},t_{0})}w\text{。}
\]
\end{lemma}

\begin{proof}
考虑函数
\[
v(s,t)=w(s_{0}+s,t_{0}+t)+\frac{b}{4}(s^{2}+t^{2})
\]
它满足
\[
\Delta v(s,t)=\Delta w(s_{0}+s,t_{0}+t)+b\ge 0\text{。}
\]
我们只需把平均值不等式\footnote{平均值不等式的这个变体在\cref{thm:appC-mean-value-inequality}中有回顾。}应用于 $v$：
\[
\begin{aligned}
w(s_{0},t_{0})=v(0)&\le\frac{1}{\pi r^{2}}\int_{B_{r}(0)}v\\
&=\frac{1}{\pi r^{2}}\int_{B_{r}(s_{0},t_{0})}w+\frac{1}{\pi r^{2}}\int_{B_{r}(0)}\frac{b(s^{2}+t^{2})}{4}\,ds\,dt\\
&=\frac{1}{\pi r^{2}}\int_{B_{r}(s_{0},t_{0})}w+\frac{br^{2}}{8}\text{，}
\end{aligned}
\]
这就证明了这条引理。
\end{proof}

下一步在 [51] 中被称为 \emph{\trHeinz}技巧（这个证明正是受了它的启发）。这个技巧在于考虑函数
\[
\begin{aligned}
f\colon[0,1]&\longrightarrow\RR\\
\rho&\longmapsto(1-\rho)^{2}\sup_{B_{\rho}(s_{0},t_{0})}w
\end{aligned}
\]
它满足 $f(0)=w(s_{0},t_{0})$ 与 $f(1)=0$。于是在 $[0,1]$ 中存在一个 $\rho_{\star}$，$f$ 在该处达到它的最大值，
\[
\max_{\rho\in[0,1]}f=f(\rho_{\star})\text{。}
\]
令
\[
c=\sup_{B_{\rho_{\star}}(s_{0},t_{0})}w=w(z_{\star})\quad\text{对 }z_{\star}\in B_{\rho_{\star}}(s_{0},t_{0})\text{，}\qquad\text{以及 }\eta=\frac{1-\rho_{\star}}{2}\text{。}
\]
注意 $\eta\in\,]0,\frac{1}{2}]$ 并且 $\rho_{\star}+\eta\le 1$，于是
\[
B_{\eta}(z_{\star})\subset B_{\rho_{\star}+\eta}(s_{0},t_{0})\subset B_{1}(s_{0},t_{0})\text{。}
\]
我们现在有\footnote{译注：原书此处第二等式右端的分母印作 $(1-\rho)^{2}$；$\rho$ 在此无定义，按前一等式（$1-\rho_{\star}-\eta=\frac{1-\rho_{\star}}{2}$）疑应为 $(1-\rho_{\star})^{2}$，此处径改。——译者注}
\[
\sup_{B_{\eta}(z_{\star})}w\le\sup_{B_{\rho_{\star}+\eta}(s_{0},t_{0})}w=\frac{f(\rho_{\star}+\eta)}{(1-\rho_{\star}-\eta)^{2}}=\frac{4f(\rho_{\star}+\eta)}{(1-\rho_{\star})^{2}}\le\frac{4f(\rho_{\star})}{(1-\rho_{\star})^{2}}=4c\text{。}
\]
因而在球 $B_{\eta}(z_{\star})$ 上我们有
\[
\Delta w\ge-aw\ge-4ac\text{。}
\]
对 $r\le\eta$ 使用\cref{thm:mean-inequality-ball}，我们便得到
\[
c=w(z_{\star})\le\frac{4ac}{8}\,r^{2}+\frac{1}{\pi r^{2}}\int_{B_{\eta}(z_{\star})}w\le\frac{ac}{2}\,r^{2}+\frac{1}{\pi r^{2}}\int_{B_{1}(s_{0},t_{0})}w\text{。}
\]
为了结束证明（这是我们的第三步，也是最后一步），令
\[
r=\frac{\eta}{\sqrt{a}}<\eta
\]
并利用最后一个不等式，我们得到
\[
c\le\frac{ac\,\eta^{2}}{2a}+\frac{a}{\pi\eta^{2}}\int_{B_{1}(s_{0},t_{0})}w\le\frac{c}{2}+\frac{a}{\pi\eta^{2}}\int_{B_{1}(s_{0},t_{0})}w\qquad\text{因为 }\eta<1\text{。}
\]
因此
\[
\frac{c}{2}\le\frac{a}{\pi\eta^{2}}\int_{B_{1}(s_{0},t_{0})}w\qquad\text{亦即 }c\eta^{2}\le\frac{2a}{\pi}\int_{B_{1}(s_{0},t_{0})}w\text{。}
\]
由此推出
\[
w(s_{0},t_{0})=f(0)\le f(\rho_{\star})=(1-\rho_{\star})^{2}c=4\eta^{2}c\le\frac{8a}{\pi}\int_{B_{1}(s_{0},t_{0})}w\text{，}
\]
这正是我们所要寻找的结论。
\end{proof}
